% 数学ノート集・共通テンプレート v2（積分論2026年度第1部の版面を参考）
% A5判。章内で定義・定理・例・注意等の番号を共通化。
% コンパイル： platex -interaction=nonstopmode math_note_template_v2.tex (2回)
%              dvipdfmx math_note_template_v2.dvi
\documentclass[a5paper,10pt,onecolumn]{jsarticle}
\usepackage[dvipdfmx,top=22mm,bottom=20mm,inner=19mm,outer=18mm,headheight=13pt,headsep=6mm]{geometry}
\usepackage{amsmath,amssymb,amsthm,mathtools}
\usepackage{enumitem,fancyhdr,titlesec}
\usepackage[dvipdfmx,bookmarks=true,colorlinks=false]{hyperref}
\usepackage{pxjahyper}

% ----- 差し替える文書情報 -----
\newcommand{\NoteTitle}{Morse理論と多様体の位相}
\newcommand{\NoteSeries}{数学ノート集}

% ----- 文字・組版 -----
% 参考PDFの質素な明朝＋Computer Modernの印象を引き継ぎながら，
% A5判で窮屈にならないよう本文の行間を少し確保する．
\linespread{1.13}
\setlength{\parindent}{1zw}
\setlength{\parskip}{0pt}
\setlength{\abovedisplayskip}{6pt plus 2pt minus 2pt}
\setlength{\belowdisplayskip}{6pt plus 2pt minus 2pt}
\setlength{\abovedisplayshortskip}{3pt plus 2pt minus 1pt}
\setlength{\belowdisplayshortskip}{4pt plus 2pt minus 1pt}
\setlist[enumerate]{topsep=3pt,itemsep=1pt,parsep=0pt,leftmargin=2.4zw}
\titleformat{\section}{\normalfont\normalsize\bfseries}{\thesection}{0.8zw}{}
\titlespacing*{\section}{0pt}{2.0ex plus .4ex}{.9ex}
\titleformat{\subsection}{\normalfont\normalsize\bfseries}{\thesubsection}{0.7zw}{}
\titlespacing*{\subsection}{0pt}{1.5ex plus .3ex}{.6ex}
\raggedbottom

% ----- 柱とページ番号：ページ上部、罫線なし -----
\fancypagestyle{mathnote}{%
  \fancyhf{}%
  \fancyhead[L]{\footnotesize \NoteSeries\ （\NoteTitle）}%
  \fancyhead[R]{\footnotesize \thepage}%
  \renewcommand{\headrulewidth}{0pt}%
  \renewcommand{\footrulewidth}{0pt}%
}
\pagestyle{mathnote}

% ----- 節内通し番号の定義・定理群 -----
\newtheoremstyle{mathnotetheorem}
 {4.5pt}{4.5pt}{\normalfont}{}%
 {\bfseries}{}{0.65zw}%
 {\thmname{#1}\thmnumber{ #2}\thmnote{（#3）}}
\theoremstyle{mathnotetheorem}
\newtheorem{theorem}{定理}[section]
\newtheorem{definition}[theorem]{定義}
\newtheorem{lemma}[theorem]{補題}
\newtheorem{proposition}[theorem]{命題}
\newtheorem{corollary}[theorem]{系}
\newtheorem{example}[theorem]{例}
\newtheorem{remark}[theorem]{注意}
\renewcommand{\proofname}{証明}
\renewcommand{\qedsymbol}{$\square$}
\makeatletter
\renewenvironment{proof}[1][\proofname]{%
  \par\pushQED{\qed}\normalfont\topsep4\p@\@plus2\p@\relax
  \trivlist\item[\hskip\labelsep\bfseries #1]\ignorespaces
}{\popQED\endtrivlist\@endpefalse}
\makeatother

% ----- 使うノートで適宜拡張する記号 -----
\newcommand{\R}{\mathbb{R}}
\newcommand{\norm}[1]{\lVert #1\rVert}
\newcommand{\abs}[1]{\lvert #1\rvert}

\begin{document}
\section*{読了後にできること}
本稿を読むための前提は，Morse関数，Hessian，Morse指数．
本稿の具体的な到達課題は，具体的な閉多様体のMorse関数を計算してEuler標数と照合する．


\section{Morse関数と臨界点}
本稿では滑らかな多様体 $M$ 上の滑らかな実関数
$f:M\to\mathbb R$ を考える．
\begin{definition}
$p$ が臨界点であるとは $df_p=0$ となることをいう．
臨界点 $p$ でのHessianを
\[
 \operatorname{Hess}_p(f)(u,v)
 =\sum_{i,j}u^iv^j\frac{\partial^2f}{\partial x^i\partial x^j}(p)
\]
で定める．これが非退化なら非退化臨界点という．
全ての臨界点が非退化である関数をMorse関数という．
Hessianの負固有値の個数を臨界点の指数という．
\end{definition}
\begin{lemma}[Hessianの座標不変性]
臨界点 $p$ でのHessianは座標変換で双線形形式として変換される．
従って非退化性と指数は座標に依存しない．
\end{lemma}
\begin{proof}
$y$ 座標への変換で二階微分の連鎖律を適用すると，
二階の座標変換の導関数には $df_p$ が係数として掛かる．
臨界点では $df_p=0$ なのでその項は消え，
Hessian行列は $J^THJ$ に変換される．
Sylvesterの慣性法則から正負の指数は不変である．
\end{proof}
\section{Morseの補題}
\begin{theorem}[Morseの補題]\label{thm:morselemma}
$p$ を指数 $\lambda$ の非退化臨界点とする．
$p$ の局所座標 $y^1,\dots,y^n$ で
\[
 f(y)=f(p)-\sum_{i=1}^{\lambda}(y^i)^2
 +\sum_{i=\lambda+1}^n(y^i)^2
\]
となるものがある．
\end{theorem}
\begin{proof}
$p$ を原点とする座標を選ぶ．
$f(0)=c$，$df_0=0$ である．
Taylorの積分剰余公式により
\[
 f(x)-c=\sum_{i,j}x^ix^jA_{ij}(x),\qquad
 A_{ij}(x)=\int_0^1(1-t)\partial_i\partial_jf(tx)\,dt.
\]
$A(x)$ は滑らかな対称行列値写像で，
$A(0)=\tfrac12\operatorname{Hess}_0(f)$ は非退化である．
$A(0)$ に対して線形座標変換を施して
$A(0)=J=\operatorname{diag}(-I_\lambda,I_{n-\lambda})$
としてよい．
$A(x)$ が $J$ に十分近い近傍で，
行列値写像 $B(x)$ を滑らかに選び
$A(x)=B(x)^T J B(x)$，$B(0)=I$ とできる．
具体的には，負部分空間と正部分空間に沿って
$A(x)$ に関するGram--Schmidt直交化を行う．
各段階の負／正の平方根の分母は $x=0$ で非零なので
近傍で滑らかに選べる．
$y=B(x)x$ と置くと
$Dy_0=B(0)=I$ だから逆写像定理によって
$y$ は局所座標となる．
\[
 f(x)-c=x^TA(x)x
 =x^TB(x)^TJB(x)x=y^TJy
\]
となり結論を得る．
\end{proof}
\begin{corollary}[臨界点の孤立性]
Morse関数の臨界点は孤立している．
コンパクト多様体上のMorse関数は有限個の臨界点を持つ．
\end{corollary}
\begin{proof}
Morse座標で $df=0$ となるのは $y=0$ のみである．
従って臨界点集合は離散的である．
また臨界点集合は $df$ の零集合なので閉であり，
コンパクト多様体ではコンパクトかつ離散の集合となる．
従って有限である．
\end{proof}
\section{臨界値のない区間}
\begin{theorem}[臨界点のない区間の同相型]\label{thm:regular}
$f^{-1}([a,b])$ がコンパクトで，
その中に臨界点がないと仮定する．
このとき劣位集合
$M^a=\{x:f(x)\le a\}$ と
$M^b=\{x:f(x)\le b\}$
は変形レトラクトによりホモトピー同値である．
\end{theorem}
\begin{proof}
$df\ne0$ の領域ではRiemann計量を固定して
\[
 X=-\frac{\nabla f}{|\nabla f|^2}
\]
と置くと $Xf=-1$ である．
コンパクトな帯 $f^{-1}([a,b])$ の近傍で
$X$ の台を切り取った滑らかなベクトル場に延長できる．
この場のフローを $\Phi_t$ とすると，
点が帯内にとどまる限り
$f(\Phi_t(x))=f(x)-t$ である．
$x\in M^b$ について
$t(x)=\max(f(x)-a,0)$ と置き
$H(x,s)=\Phi_{s\,t(x)}(x)$ とする．
$x\in M^a$ なら $t(x)=0$ なので固定され，
$s=1$ で像は $M^a$ に入る．
$H$ は連続であるから強変形レトラクトを与える．
\end{proof}
\section{臨界点を通過するときの変化}
\begin{lemma}[局所モデルの相対位相]
$Q(u,v)=-|u|^2+|v|^2$，
$u\in\mathbb R^\lambda$，
$v\in\mathbb R^{n-\lambda}$ とする．
零点近くの十分小さな閉球 $B$ について
$Q$ の劣位集合が負の臨界値を越えると，
相対的に $\lambda$ 次元の胞体を付加する操作が生じる．
\end{lemma}
\begin{proof}
局所的な対 $(Q\le\varepsilon,Q\le-\varepsilon)$ を考える．
負方向の円板
$D^\lambda\times\{0\}$ と
正方向の円板 $\{0\}\times D^{n-\lambda}$ を用い，
$Q(u,v)$ の値を増加させない方向へ $v$ を縮める
ホモトピー
$(u,v)\mapsto(u,(1-t)v)$ を考える．
これにより正方向を収縮し，残る変化は
$|u|^2\le\varepsilon$ の円板をその境界
$S^{\lambda-1}$ に沿って加える場合に帰着する．
一般の小球による切断では境界でこのホモトピーを
切り替える必要があるが，
半径方向のcutoff関数を用いれば，
球の外側を固定して同じ相対変形を得る．
\end{proof}
\begin{theorem}[Morseの胞体付加定理]\label{thm:cell}
$f^{-1}([a,b])$ がコンパクトで，
その中の臨界点がただ一つ $p$ で，
指数 $\lambda$，臨界値 $c\in(a,b)$ とする．
このとき $M^b$ は $M^a$ に
$\lambda$ 次元胞体を一つ付加した空間とホモトピー同値である．
\end{theorem}
\begin{proof}
Morseの補題によって $p$ の近傍で
$f-c=-|u|^2+|v|^2$ と表示できる．
$[a,c-\delta]$ と $[c+\delta,b]$ には臨界点がないので，
定理\ref{thm:regular}を使ってこの部分の劣位集合を
フローで変形できる．
残る臨界点近傍では前補題の局所モデルを使用する．
局所モデルの外の領域では正規化した負勾配流を用い，
境界でcutoff関数を選んで両変形を一致させると，
$M^b$ のホモトピー型は
$M^a$ と
$D^\lambda$ を境界 $S^{\lambda-1}$ に沿って貼り合わせた
空間に等しくなる．
\end{proof}
\section{Morse不等式}
\begin{theorem}[弱Morse不等式と強Morse不等式]\label{thm:inequality}
$M$ をコンパクト滑らかな多様体，$f$ をMorse関数とし，
指数 $k$ の臨界点の個数を $m_k$，
体 $\mathbb F$ 係数のBetti数を
$b_k=\dim_{\mathbb F}H_k(M;\mathbb F)$ とする．
このとき
\[
 m_k\ge b_k,\qquad
 \sum_{j=0}^q(-1)^{q-j}m_j
 \ge\sum_{j=0}^q(-1)^{q-j}b_j
\]
が全ての $q$ について成立し，
$q=n$ では等号となる．
\end{theorem}
\begin{proof}
Morse胞体付加定理によって
$M$ は指数 $k$ の臨界点ごとに $k$ 胞体を一つ持つ
有限CW複体のホモトピー型を持つ．
体係数の胞体鎖群を $C_k$ とすれば
$\dim C_k=m_k$ である．
境界写像 $\partial_k:C_k\to C_{k-1}$ の階数を
$r_k$ と置くと
\[
 b_k=\dim\ker\partial_k-\dim\operatorname{im}\partial_{k+1}
 =m_k-r_k-r_{k+1}.
\]
従って $m_k-b_k=r_k+r_{k+1}\ge0$．
また交代和を取ると
\[
 \sum_{j=0}^q(-1)^{q-j}(m_j-b_j)=r_{q+1}\ge0
\]
となる．$q=n$ では $C_{n+1}=0$ なので $r_{n+1}=0$，
従って等号が成立する．
\end{proof}
\begin{remark}
胞体付加定理の詳細な技術的部分，
とくに臨界点近傍と外側の勾配フローの貼り合わせは
最終校閲で補強する．
\end{remark}

\section{到達課題の完全解答}
\noindent\textbf{球面の高さ関数の臨界点．}
\begin{proposition}
$S^2\subset\mathbb R^3$ 上の高さ関数
$h(x,y,z)=z$ は北極と南極だけを臨界点として持ち，
南極のMorse指数は0，北極の指数は2である．
\end{proposition}
\begin{proof}
$S^2$ の接空間は $p^\perp$ であるから，
$dh_p$ が接空間上零になるのは
鉛直方向が法線 $p$ に平行なときであり，
$p=(0,0,\pm1)$ に限られる．
各極付近で $(x,y)$ を局所座標とすると
$z=\pm\sqrt{1-x^2-y^2}$．
原点でのHessianは北極で $-I_2$，
南極で $I_2$ だからいずれも非退化で，
負固有値の個数はそれぞれ2と0である．
\end{proof}
\begin{corollary}
Morse臨界点の交代和は
$1+1=2=\chi(S^2)$ となる．
\end{corollary}
\begin{proof}
指数0の臨界点が一つ，指数2の臨界点が一つ，
指数1の臨界点は存在しない．
Morse理論のEuler標数公式に代入する．
\end{proof}
\noindent\textbf{到達確認．}
$S^1$ 上の高さ関数 $h(x,y)=y$ の
Morse指数を求めよ．
\par\noindent\textbf{解答．}
最下点で指数0，最上点で指数1である．
交代和は $1-1=0=\chi(S^1)$．
\end{document}
